\chapter{Intro}\label{ch:intro}
It's reasonable to expect that quantum computers will eventually be practically useful.
However, it's also reasonable to expect that any useful quantum computer will have to be fault-tolerant.
That is, the quantum information used in the computation will have to be encoded with some redundancy into a larger quantum system, such that when faults inevitably occur on the physical components, the system as a whole can tolerate this, in that the desired computation can still be performed.

This thesis can be viewed as a combination of two things.
The first is the following mantra: every Clifford circuit with Pauli measurements is a quantum error correcting code.
This viewpoint is the essence of dynamic quantum error correction, a topic that the diligent reader will notice makes up 80\% of the title of this thesis.
The remaining 20\% -- `picturing' -- refers to the ZX-calculus, and perhaps specifically to a tool within this called Pauli webs.

%Humans have solved the problem of fault-tolerance once before: in classical computing.
%In this familiar realm, the physical components forming the building blocks of computers might experience faults at a rate

Dynamic quantum error correction -- more accurately, dynamic stabilizer codes -- burst onto the scene with Hastings and Haah's paper \textit{Dynamical Logical Qubits}~\cite{Hastings_2021}.
%In a nutshell, whereas a stabilizer code $S$ encodes $k$-qudit states into a certain fixed subspace $V_S$ of a larger $n$-qudit system $\mathcal{H}$, a dynamic stabilizer code encodes $k$-qudit states into a sequence $V_0, V_1, \ldots$ of subspaces of $\HHH$.
In a nutshell, whereas a stabilizer code $\SSS$ is implemented by repeatedly measuring a fixed set of commuting Pauli operators, which encodes information into a fixed subspace $\VVV_\SSS$ of a larger system $\mathcal{H}$, a dynamic stabilizer code is implemented by performing an arbitrary sequence of Pauli measurements and Clifford unitaries, which encodes information into a sequence $\VVV_0, \VVV_1, \ldots$ of subspaces of $\HHH$.
When Hastings and Haah's paper appeared, it was not only theoretically interesting, as it expanded our ideas around what it means to be a quantum error correcting code, but the hope was also that it might also bring practical advantages -- e.g. in circumventing no-go theorems that shackle static quantum codes.
While the former remains true, the latter did not come to pass.
Indeed, as we'll discuss in this thesis, the practical advantages of dynamic stabilizer codes are very much a trade-off rather than a free lunch.
As such, after an initial hot dynamic QEC summer, interest in the topic feels anecdotally as though it's waning, as we come to understand it better and realise it isn't as alien as it initially may have seemed.

The ZX-calculus~\cite{Coecke_2011}, on the other hand, has experienced a very different trajectory.
No bursting onto the scene here; instead, over the course of the last decade or so, the ZX-calculus has slowly won over skeptics to take its current place as a useful tool in the quantum computer scientist's toolbox.
The ZX-calculus is a graphical language for reasoning about quantum mechanics, in the same way that linear algebra is a symbolic language for reasoning about quantum mechanics.
The secret to its success (so I claim!) is simply that humans have been designed by evolution to be better at reasoning graphically than symbolically.
In addition, the ZX-calculus is a fairly simple language; there are just two generators, the green and red spiders, and just seven rules are needed to fully reason about quantum mechanics up to the second level of the Clifford hierarchy.
Happily, this fragment contains many of the QEC primitives we wish to reason about, and the first half of this thesis takes place inside it.

Within the ZX-calculus, Pauli webs~\cite{McEwen_2023,Bombin_2024} are a graphical language for reasoning about the stabilizer formalism.
Since the stabilizer formalism is the basis for defining dynamic stabilizer codes, it isn't surprising that Pauli webs have become an invaluable tool for reasoning about such codes.
That said, the stabilizer formalism is also the basis for \textit{static} stabilizer codes, and these seem to have done alright without Pauli webs for the last thirty years -- are Pauli webs really all that necessary?

The difference is that, even for static stabilizer codes $\SSS$ that have complicated algebraic definitions, once they've been defined they're pretty simple to reason about.
One repeatedly measures the same set of generators, which repeatedly projects the qubits into the fixed codespace $\VVV_\SSS$ in the absence of errors.
Errors are detected by comparing the results of these measurements in consecutive rounds, and logical operators are just elements of a particular static quotient group denoted $\NNN(\SSS)/\SSS$.
Dynamic stabilizer codes, on the other hand, can get rather fiddly rather quickly -- just from writing down the sequence of measurements and Cliffords to be performed, it's not at all obvious which measurements might be combined to detect errors or what a representative of a particular logical operator might be after a certain number of operations.
Pauli webs make all this graphical, and therefore a lot more tractable for us humans and our good old visual cortexes.

In the second half of this thesis, we essentially extend the ideas of the first half into the world of continuous variable (CV) quantum information theory.
Here we meet GKP codes, the analogue of Pauli stabilizer codes, and are introduced to the green and red spiders' Focked up cousin.
We're able to generalise Pauli webs to this new setting, and then use them to introduce and explore dynamic CV quantum error correction.
Moreover, we do this in \textit{fault-equivalent} manner~\cite{rodatz2025faulttoleranceconstruction}, guaranteeing that the graphical manipulations we're performing do not adversely impair the behaviour of our diagrams under noise.
